%====================================================================================
%  single-ref.tex  —  FORMAT HOLDER / SKELETON for single-* (one-page) resumes
%  ----------------------------------------------------------------------------------
%  NOT an application resume. Copy single-ComVis.tex (the GOLD exemplar) to build a
%  real one; use THIS file as the rules-in-context cheat sheet. See CLAUDE.md sec 3.
%
%  HARD RULES (all measured in this exact layout; usable bullet width = 537.7pt):
%   1. Exactly ONE page. Every \resumeItem must be ONE line, no wrap.
%   2. Pack each bullet full: target 108-114 chars plain prose (wraps ~115-118).
%      Bold costs ~15% width, ALL-CAPS ~42% — spend them deliberately.
%   3. >= 30 content lines total. Per section: Education 3 achv + Skills line;
%      Professional Exp >=3 entries x >=3 bullets; Research >=2 x >=3;
%      Projects >=3 x >=3; Publications 3.
%   4. One complete idea per bullet (a full sentence). Never split one sentence
%      across multiple bullets to pad the count.
%   5. ALWAYS include: Publications + Ericsson internship + ABB thesis.
%   6. Publications block = LITERAL copy-paste from single-GenMod.tex (normal font).
%   7. Order EVERY section reverse-chronologically by END date (latest first).
%   8. Dates = month ranges ("Mon YYYY -- Mon YYYY"), never bare years.
%   9. No em dashes / no dash-as-separator. Titles use ":"; GPA in "(...)";
%      "--" only for date ranges.
%  10. Favor concepts + results over obscure names. STRIP proprietary/research
%      tool/model/dataset/benchmark names (COLMAP, OpenMVS, ROUGE, ARKitScenes,
%      FairMOT, ...) and raw math. KEEP field-standard ATS keywords (SIFT, RANSAC,
%      CLIP, Hough, ResNet, ViT) and a name only if central AND explained inline.
%  11. To fit one page, tighten vertical \vspace only (never margins/font).
%  12. Verify: pages==1, 0 overfull, >=30 "^•" markers, each publication one line.
%====================================================================================

\documentclass[letterpaper, 11pt]{article}
\usepackage{latexsym}
\usepackage[empty]{fullpage}
\usepackage{titlesec}
\usepackage{marvosym}
\usepackage[usenames, dvipsnames]{color}
\usepackage{verbatim}
\usepackage{enumitem}
\usepackage[hidelinks]{hyperref}
\usepackage{fancyhdr}
\usepackage[english]{babel}
\usepackage{tabularx}
\usepackage{fontawesome5}
\usepackage{multicol}
\setlength{\multicolsep}{-3.0pt}
\setlength{\columnsep}{-1pt}
\input{glyphtounicode}
\usepackage{vwcol}
\usepackage{comment}
\usepackage{hyperref}
\hypersetup{
	colorlinks=true,
	linkcolor=blue,
	filecolor=black,
	citecolor=black,
	urlcolor=Blue,
}

\pagestyle{fancy}
\fancyhf{} % clear all header and footer fields
\fancyfoot{}
\renewcommand{\headrulewidth}{0pt}
\renewcommand{\footrulewidth}{0pt}

% Adjust margins
\addtolength{\oddsidemargin}{-0.5in}
\addtolength{\evensidemargin}{-0.5in}
\addtolength{\textwidth}{1.19in}
\addtolength{\topmargin}{-.5in}
\addtolength{\textheight}{1.2in}

\urlstyle{same}

\raggedbottom
\raggedright
\setlength{\tabcolsep}{0in}

% Sections formatting
\titleformat{\section}{ \vspace{-4pt}\scshape\raggedright\large\bfseries }{}{0em}{}[
\color{Blue}{}
\titlerule
\vspace{-5pt}
]

% Ensure that generate pdf is machine readable/ATS parsable
\pdfgentounicode=1

%-------------------------
% Custom commands
\newcommand{\resumeItem}[1]{ \item\small{ {#1 \vspace{-2pt}} } }

\newcommand{\resumeSubheading}[4]{
\vspace{-2pt}
\item
\begin{tabular*}{1.0\textwidth}[t]{l@{\extracolsep{\fill}}r}
	\textbf{#1}       & {\small #2} \\
     [-2pt]
	\textit{{\small#3}} & \textit{\small #4} \\
\end{tabular*}
\vspace{-10pt}
}

\newcommand{\resumeProjectHeading}[2]{
\vspace{-2pt}
\item
\begin{tabular*}{1.0\textwidth}[t]{l@{\extracolsep{\fill}}r}
	\textbf{#1}  & {\small #2} \\
\end{tabular*}
\vspace{-18pt}
}

\newcommand{\resumeSubItem}[1]{\resumeItem{#1}
\vspace{-4pt}}

\renewcommand{\labelitemi}{$\vcenter{\hbox{\tiny$\bullet$}}$}
\renewcommand{\labelitemii}{$\vcenter{\hbox{\tiny$\bullet$}}$}

\newcommand{\resumeSubHeadingListStart}{\begin{itemize}[leftmargin=0.0in, label={}]}
\newcommand{\resumeSubHeadingListEnd}{\end{itemize}\vspace{-1em}}
\newcommand{\resumeItemListStart}{\vspace{-0.3em}\begin{itemize}[leftmargin=0.25in]}
\newcommand{\resumeItemListEnd}{\end{itemize}\vspace{-7.5pt}}  % <- tighten THIS (more negative) if a real resume spills to page 2

%-------------------------------------------
%%%%%%  RESUME STARTS HERE  %%%%%%%%%%%%%%%%%%%%%%%%%%%%

\begin{document}

%----------HEADING---------- (FIXED boilerplate — reuse verbatim; links: Phone, Email, LinkedIn, GitHub, Portfolio, Location)
\begin{center}
    {\Huge \scshape Gawtam Chithra Ramesh} \\
    \vspace{3pt}
    \small
    { \href{tel:+46764518589}{{\faMobile*} \textcolor{black}{(+46) 76541-8589}} ~
    \href{mailto:gawtamcr3@gmail.com}{{\faPaperPlane} \textcolor{black}{gawtamcr3@gmail.com}} ~
    \href{https://linkedin.com/in/gawtamcr/}{{\faLinkedinIn} \textcolor{black}{Gawtam C R}} ~
    \href{https://github.com/gawtamcr}{{\faGithub} \textcolor{black}{gawtamcr}} ~
    \href{https://gawtamcr.github.io/}{{\faGlobe} \textcolor{black}{Portfolio}} ~
    \href{https://maps.google.com/?q=Stockholm,Sweden}{{\faMapMarker*} \textcolor{black}{Stockholm, Sweden}}
    }
    \vspace{-8pt}
\end{center}

%-----------EDUCATION----------- (FIXED boilerplate; GPA in parentheses, NOT after a dash. Latest end date first.)
\section{\color{Blue}Education}
    \resumeSubHeadingListStart
    \resumeSubheading{KTH Royal Institute of Technology}{Aug 2024 -- Jun 2026}{Master of Science in Systems, Controls, and Robotics \textnormal{(GPA 4.62/5.0)}}{Stockholm, Sweden}
    \resumeSubheading{Indian Institute of Technology Madras}{Aug 2019 -- Jul 2024}{Dual Degree (B.Tech Engineering Design and M.Tech Robotics) \textnormal{(GPA 7.9/10)}}{Chennai, India}
    \vspace{-0.3em}
    \small{
    \begin{itemize}[leftmargin=0.15in]
        \setlength{\itemsep}{-0.3em}
        \item Awarded (among 30 of 800+) the IIT Madras \textbf{Young Research Fellowship} based on remarkable research potential.
        \item Secured \textbf{97.64} Percentile in JEE Advanced and \textbf{99.10} Percentile in JEE Mains among \textbf{1.5M} candidates.
        \item Secured an All India Rank of \textbf{52} out of 50k in Undergraduate Common Entrance Examination for Design.
    \end{itemize}
    }
    \vspace{-0.7em}
    \textbf{Skills} : [Tailor to the role; ~12-13 keywords, comma-separated, ONE line. Lead with the job's must-haves.]
    \resumeSubHeadingListEnd

%-----------EXPERIENCE----------- (>=3 entries, >=3 bullets each; reverse-chron by end date; ABB + Ericsson are MANDATORY)
\section{\color{Blue}Professional Experience}
\resumeSubHeadingListStart

    % --- WORKED EXAMPLE (mandatory entry). Note: title uses ":" not a dash; bullets ~108-114 chars, concept+mechanism+result.
    \resumeSubheading{Master's Thesis: [Short Thesis Title, no em dash]}{Jan 2026 -- Jun 2026}{ABB Robotics}{V\"{a}ster\aa s, Sweden}
    \resumeItemListStart
        \resumeItem{Designed a generative flow-matching planner guided at inference by Signal Temporal Logic robustness gradients.}
        \resumeItem{Steered the pretrained flow model toward new temporal-logic specifications at inference without retraining.}
        \resumeItem{Benchmarked on long-horizon reach-avoid navigation against state-of-the-art specification-guided motion planners.}
    \resumeItemListEnd

    % --- WORKED EXAMPLE (mandatory entry). "Hydra" kept because it is central AND explained inline.
    \resumeSubheading{Device Research Intern: 3D Scene Graphs}{Jul 2025 -- Dec 2025}{Ericsson Research}{Lund, Sweden}
    \resumeItemListStart
        \resumeItem{Integrated instance segmentation and tracking into Hydra, a real-time 3D scene-graph construction framework.}
        \resumeItem{Built a synchronized RGB-D inpainting pipeline that removes privacy-sensitive objects from scene graphs online.}
        \resumeItem{Benchmarked state-of-the-art inpainting models to balance reconstruction fidelity against real-time latency.}
    \resumeItemListEnd

    % --- HINT TEMPLATE (replace with the most role-relevant 3rd experience; delete brackets).
    \resumeSubheading{[Role Title: concise, no em dash]}{[Mon YYYY -- Mon YYYY]}{[Company]}{[City, Country]}
    \resumeItemListStart
        \resumeItem{[Action verb + what you built/did + the mechanism + the outcome. One sentence, ~108-114 chars, no obscure names.]}
        \resumeItem{[Second distinct accomplishment with a concrete result or metric if you have one (e.g. 86.2\% accuracy).]}
        \resumeItem{[Third accomplishment. Keep it a full sentence; field-standard keywords OK, niche tool/dataset names out.]}
    \resumeItemListEnd

\resumeSubHeadingListEnd

%-----------RESEARCH----------- (>=2 entries, >=3 bullets each; reverse-chron by end date)
\section{\color{Blue}Research Experience}
\resumeSubHeadingListStart

    % --- HINT TEMPLATE. Third field = advisor + lab/department (optionally \href-linked) + institution.
    \resumeSubheading{[Project Title]}{[Mon YYYY -- Mon YYYY]}{\href{[url]}{[Prof. Name]}, [Lab / Department], [Institution]}{[City, Country]}
    \resumeItemListStart
        \resumeItem{[What problem you tackled and the approach — one full technical sentence packed to the line width.]}
        \resumeItem{[The method/framework you built and what it produces.]}
        \resumeItem{[Evaluation + result. Use generic metric names ("text-similarity metrics") not niche ones (ROUGE/BERTScore).]}
    \resumeItemListEnd

    \resumeSubheading{[Project Title]}{[Mon YYYY -- Mon YYYY]}{\href{[url]}{[Prof. Name]}, [Department], [Institution]}{[City, Country]}
    \resumeItemListStart
        \resumeItem{[Concept-first sentence: what was reconstructed/built/estimated, via which standard technique.]}
        \resumeItem{[Next step in the pipeline and its output/integration.]}
        \resumeItem{[Final stage / downstream use. Full sentence, one line.]}
    \resumeItemListEnd

\resumeSubHeadingListEnd

%-----------PROJECTS----------- (>=3 entries, >=3 bullets each; reverse-chron by end date)
\section{\color{Blue}Technical Projects}
\resumeSubHeadingListStart

    % Heading 2nd arg is the date range; 1st arg is "[Project Name] | [Course Code, Institution]".
    \resumeProjectHeading{[Project Name] $|$ \textnormal{\small [Course Code, Institution]}}{[Mon YYYY -- Mon YYYY]}
    \resumeItemListStart
        \resumeItem{[What you implemented and how — full sentence, ~108-114 chars, field-standard keywords welcome.]}
        \resumeItem{[Second component or design choice.]}
        \resumeItem{[Result/metric (e.g. "reaching 86.2\% validation accuracy and 0.917 ROC-AUC").]}
    \resumeItemListEnd

    \resumeProjectHeading{[Project Name] $|$ \textnormal{\small [Course Code, Institution]}}{[Mon YYYY -- Mon YYYY]}
    \resumeItemListStart
        \resumeItem{[Bullet 1 — one idea, one line.]}
        \resumeItem{[Bullet 2 — one idea, one line.]}
        \resumeItem{[Bullet 3 — one idea, one line.]}
    \resumeItemListEnd

    \resumeProjectHeading{[Project Name] $|$ \textnormal{\small [Course Code, Institution]}}{[Mon YYYY -- Mon YYYY]}
    \resumeItemListStart
        \resumeItem{[Bullet 1 — one idea, one line.]}
        \resumeItem{[Bullet 2 — one idea, one line.]}
        \resumeItem{[Bullet 3 — one idea, one line.]}
    \resumeItemListEnd

\resumeSubHeadingListEnd

%-----------PUBLICATIONS----------- (RULE 6: paste this block VERBATIM from single-GenMod.tex; normal font, NOT \small)
\section{\color{Blue}Lead Author Publications}
\begin{itemize}[leftmargin=0.15in]
    \setlength{\itemsep}{-5pt}
    \item Robust Object-layer Construction of 3DSG Using Instance Segmentation (IMPROVE26)(\textbf{Best Paper Award})
    \item Synchronized RGB-D Inpainting for Privacy-Aware 3D Scene Graph Construction.(\textbf{ICPR26}).(Under Review)
    \item Scene Understanding in Deformable Object Manipulation via Taxonomy-Guided VLMs.(\textbf{IROS25} Workshop)
\end{itemize}

\end{document}
